\Needspace{12\baselineskip}
\begin{center}
\textit{Forms with determinant 0.}
\end{center}

\Needspace{8\baselineskip}
\Article{215}
So far we have excluded forms with determinant 0 from all our investigations. We must add a few things about them to make our theory complete in every respect. We have proved generally that if a form with determinant $D$ contains a form with determinant $D'$, then $D'$ is a multiple of $D$. So a form with determinant $=0$ clearly cannot contain any form whose determinant is not also $=0$. Only two problems therefore remain: $1^{\mathrm{o}}$ \textit{given two forms $f,F$, the latter with determinant 0, decide whether the former contains the latter, and if it does, give all transformations of the former into the latter;} $2^{\mathrm{o}}$ \textit{find all representations of a given number by a given form with determinant 0.} The first problem requires one
\SourceLeaf{217}
method when the determinant of the first form $f$ is also 0, and another when it is not 0. We will now set all this out.

I. First we observe that every form $axx+2bxy+cyy$ with determinant $bb-ac=0$ can be written as $m(gx+hy)^2$, where $g,h$ are relatively prime and $m$ is an integer. For let $m$ be the greatest common divisor of $a,c$, taken with the same sign as these numbers themselves (we easily see that they cannot have opposite signs). Then $\frac{a}{m},\frac{c}{m}$ are relatively prime nonnegative integers; their product is $=\frac{bb}{mm}$, \textit{i.e.}, a square, so they themselves are squares (Article~21). Let $\frac{a}{m}=gg$, $\frac{c}{m}=hh$. Then $g,h$ are also relatively prime, $gghh=\frac{bb}{mm}$, and $gh=\pm\frac{b}{m}$. Clearly, then,
\[
m(gx\pm hy)^2=axx+2bxy+cyy
\]

Now let two forms $f,F$ be given, both with determinant 0, and let
\[
f=m(gx+hy)^2,\quad F=M(GX+HY)^2
\]
where $g$ is relatively prime to $h$, and $G$ to $H$. I claim that if $f$ contains $F$, then $m$ either equals $M$ or at least divides $M$, and the quotient is a square. Conversely, if $\frac{M}{m}$ is an integer square, $F$ is contained in $f$. For suppose $f$ passes by the substitution
\[
x=\alpha X+\beta Y,\quad y=\gamma X+\delta Y
\]
into $F$. Then
\[
\frac{M}{m}(GX+HY)^2=((\alpha g+\gamma h)X+(\beta g+\delta h)Y)^2
\]
from which we easily get that $\frac{M}{m}$ is a square. Put it $=ee$. Then
\[
e(GX+HY)=\pm((\alpha g+\gamma h)X+(\beta g+\delta h)Y),\quad\text{\textit{i.e.}}
\]
\[
\pm eG=\alpha g+\gamma h,\quad\pm eH=\beta g+\delta h
\]
So if $\mathfrak{G},\mathfrak{H}$ are chosen so that $\mathfrak{G}G+\mathfrak{H}H=1$, we have
\[
\pm e=\mathfrak{G}(\alpha g+\gamma h)+\mathfrak{H}(\beta g+\delta h),\quad\text{so an \textit{integer.}\quad\textit{Q.E.D.}}
\]
Conversely, suppose $\frac{M}{m}$ is an integer square $=ee$. Then $f$ contains $F$. Namely, we can choose integers $\alpha,\beta,\gamma,\delta$ so that
\SourceLeaf{218}
\[
\alpha g+\gamma h=\pm eG,\quad \beta g+\delta h=\pm eH
\]
For choose integers $\mathfrak{g},\mathfrak{h}$ so that $\mathfrak{g}g+\mathfrak{h}h=1$. We satisfy these equations by putting
\[
\begin{aligned}
\alpha&=\pm eG\mathfrak{g}+hz, &\quad\gamma&=\pm eG\mathfrak{h}-gz\\
\beta&=\pm eH\mathfrak{g}+hz', &\quad\delta&=\pm eH\mathfrak{h}-gz'
\end{aligned}
\]
whatever integer values we give $z,z'$. So $F$ is contained in $f$. \textit{This proves the claim.} We also easily see that these formulas give all possible values of $\alpha,\beta,\gamma,\delta$, \textit{i.e.}, all transformations of $f$ into $F$, if $z,z'$ range over all integers.

II. Given two forms $f=axx+2bxy+cyy$ with determinant not $=0$, and $F=M(GX+HY)^2$ with determinant $=0$ (where, as before, $G,H$ are relatively prime), I claim \textit{first} that if $f$ contains $F$, then $M$ can be represented by $f$; \textit{second}, if $M$ can be represented by $f$, then $F$ is contained in $f$; \textit{third}, if in this case all representations of $M$ by $f$ are given by $x=\xi,y=\upsilon$, then all transformations of $f$ into $F$ are given by $G\xi,H\zeta,G\upsilon,H\upsilon$. We prove all this as follows.

$1^{\mathrm{o}}$ Suppose $f$ passes into $F$ by the substitution $\alpha,\beta,\gamma,\delta$, and choose $\mathfrak{G},\mathfrak{H}$ so that $\mathfrak{G}G+\mathfrak{H}H=1$. Clearly, putting $x=\alpha\mathfrak{G}+\beta\mathfrak{H}$, $y=\gamma\mathfrak{G}+\delta\mathfrak{H}$ makes the value of $f$ equal to $M$, so $M$ is representable by $f$.

$2^{\mathrm{o}}$ Suppose $a\xi\xi+2b\xi\upsilon+c\upsilon\upsilon=M$. Clearly, by the substitution $G\xi,H\xi,G\upsilon,H\upsilon$, the form $f$ passes into $F$. To see that

$3^{\mathrm{o}}$ in this case the substitution $G\xi,H\zeta,G\upsilon,H\upsilon$ gives \textit{all} transformations of $f$ into $F$, if $\xi,\upsilon$ give all values of $x,y$ that make $f=M$, we argue as follows. Let $\alpha,\beta,\gamma,\delta$ be any transformation of $f$ into $F$, and, as before, $\mathfrak{G}G+\mathfrak{H}H=1$. Among the values of $x,y$ are then
\[
x=\alpha\mathfrak{G}+\beta\mathfrak{H},\quad y=\gamma\mathfrak{G}+\delta\mathfrak{H}
\]
which give the substitution
\SourceLeaf{219}
\[
G(\alpha\mathfrak{G}+\beta\mathfrak{H}),\quad H(\alpha\mathfrak{G}+\beta\mathfrak{H}),\quad G(\gamma\mathfrak{G}+\delta\mathfrak{H}),\quad H(\gamma\mathfrak{G}+\delta\mathfrak{H})
\]
or
\[
\begin{aligned}
&\alpha+\mathfrak{H}(\beta G-\alpha H), &\quad&\beta+\mathfrak{G}(\alpha H-\beta G)\\
&\gamma+\mathfrak{H}(\delta G-\gamma H), &&\delta+\mathfrak{G}(\gamma H-\delta G)
\end{aligned}
\]
But since
\[
\begin{aligned}a(\alpha X+\beta Y)^2+2b(\alpha X+\beta Y)(\gamma X+\delta Y)&{}\\
{}+c(\gamma X+\delta Y)^2&=M(GX+HY)^2\end{aligned}
\]
we have
\[
\begin{aligned}
a(\alpha\delta-\beta\gamma)^2&=M(\delta G-\gamma H)^2\\
c(\beta\gamma-\alpha\delta)^2&=M(\beta G-\alpha H)^2
\end{aligned}
\]
and therefore (since the determinant of $f$, multiplied by $(\alpha\delta-\beta\gamma)^2$, equals the determinant of $F$, \textit{i.e.}, $=0$, so that also $\alpha\delta-\beta\gamma=0$),
\[
\delta G-\gamma H=0,\quad \beta G-\alpha H=0
\]
The substitution thus becomes $\alpha,\beta,\gamma,\delta$, which shows that the formula given supplies \textit{all} transformations of $f$ into $F$.

III. It remains to show how to give all representations of a given number by a given form with determinant 0. Let the form be $m(gx+hy)^2$. Clearly the number must be divisible by $m$, and the quotient must be a square. If we put the given number $=mee$, then for the same values of $x,y$ that make $m(gx+hy)^2=mee$, clearly $gx+hy$ is either $=+e$ or $=-e$. So we obtain all representations by finding all integer solutions of the linear equations $gx+hy=e$, $gx+hy=-e$. These are known to be solvable, since $g,h$ are relatively prime, as assumed. Namely, if we choose $\mathfrak{g},\mathfrak{h}$ so that $\mathfrak{g}g+\mathfrak{h}h=1$, we satisfy the first equation by putting $x=\mathfrak{g}e+hz$, $y=\mathfrak{h}e-gz$, and the second by putting $x=-\mathfrak{g}e+hz$, $y=-\mathfrak{h}e-gz$, where $z$ is any integer. These formulas give \textit{all} integer values of $x,y$ if $z$ ranges over all integers.

